\documentclass[11pt]{article}

\usepackage[margin=1in]{geometry}
\usepackage{amsmath}
\usepackage{amssymb}
\usepackage{graphicx}
\usepackage{booktabs}
\usepackage{array}
\usepackage[hidelinks]{hyperref}

\title{Locating Executable Files in an Extracted Cyberpunk~2077 Installation:\\ A Recursive Glob Audit of the Game Tree and Redist Folder}
\author{[Author placeholder]\\ \texttt{[author@example.com]}}
\date{September 24, 2026}

% TODO: Replace the author placeholder with the actual author name and affiliation.
% TODO: The source session cited no external literature; the References section
% therefore contains only an internal pointer to the prior session work. Do not
% add external citations without grounding them in the session summary.

\begin{document}

\maketitle

\begin{abstract}
This paper documents a systematic audit of executable (``.exe'') files within
an extracted Cyberpunk~2077 game installation located at
\texttt{/Volumes/One Touch/mdata/games/cbp/Cyberpunk 2077/}. The installation
was previously extracted from a zstd-compressed ZIP archive in prior session
work. A recursive glob search with the pattern \texttt{**/*.exe} was executed
over the game tree and, separately, over the \texttt{Redist} folder. The audit
identified 12 executables in total: 7 game-related binaries within the game
tree and 5 redistributable runtime installers within \texttt{Redist/}. The
observed layout is consistent with a standard pre-installed Cyberpunk~2077
repack, with game binaries under \texttt{bin/x64}, modding tools under
\texttt{tools/redmod/bin}, and redistributable runtimes under
\texttt{Redist/}. The primary launcher of interest is
\texttt{Cyberpunk2077.exe}, with \texttt{Run Me!.bat} available for the setup
steps.
\end{abstract}

\section{Introduction}
\label{sec:introduction}

Pre-installed (``repacked'') game distributions are typically delivered as a
single archive that, once extracted, must be verified before use: the user
needs to confirm that the installation is complete, identify which executable
launches the game, and distinguish game binaries from the runtime dependencies
that merely support them. After a large game archive has been extracted, a
first-order integrity check is to enumerate the executable files present in the
installation tree and to map them to their roles.

This paper reports such an audit for a Cyberpunk~2077 installation that was
extracted from a zstd-compressed ZIP archive in prior session work (see
Section~\ref{sec:background}). The extracted game tree resides at
\texttt{/Volumes/One Touch/mdata/games/cbp/Cyberpunk 2077/}. The session that
produced this paper pursued the following research questions:

\begin{enumerate}
  \item \textbf{RQ1.} Which executable (``.exe'') files are present in the
        extracted installation, and where are they located?
  \item \textbf{RQ2.} Which of these executables are game binaries, and which
        are redistributable runtime dependencies?
  \item \textbf{RQ3.} Which executable serves as the primary launcher of the
        game?
\end{enumerate}

The remainder of this paper is organized as follows.
Section~\ref{sec:background} summarizes the prior extraction work that produced
the installation under audit. Section~\ref{sec:methods} describes the search
procedure. Section~\ref{sec:results} reports the executables found, and
Section~\ref{sec:discussion} interprets the resulting layout.
Section~\ref{sec:conclusion} concludes.

\section{Background}
\label{sec:background}

The installation under audit was produced in prior session work, in which the
Cyberpunk~2077 game archive was extracted from a zstd-compressed ZIP archive.
The extraction target was
\texttt{/Volumes/One Touch/mdata/games/cbp/Cyberpunk 2077/} on an external
drive. The present session is a follow-up: it does not re-examine the
extraction itself, but audits the resulting installation tree to locate all
executable files and to identify the launcher.

\section{Methods}
\label{sec:methods}

\subsection{Search procedure}
\label{sec:procedure}

Executable files were located by recursive glob search using the pattern
\texttt{**/*.exe}, which matches every file whose name ends in
\texttt{.exe} at any depth of the directory tree. Two searches were performed:

\begin{itemize}
  \item \textbf{Game tree.} A recursive glob search over the full game tree at
        \texttt{/Volumes/One Touch/mdata/games/cbp/Cyberpunk 2077/}.
  \item \textbf{Redist folder.} A separate recursive glob search over the
        \texttt{Redist} folder, which in pre-installed distributions typically
        holds redistributable runtime installers rather than game binaries.
\end{itemize}

\subsection{Scope and limitations}
\label{sec:scope}

The audit is limited to filename-based enumeration: it establishes the
presence and location of executable files but performs no binary inspection,
no integrity verification, and no execution of any of the files found. The
classification of executables into roles (game binary, launcher, modding tool,
runtime dependency) is based on their paths and names within the standard
installation layout.

\section{Results}
\label{sec:results}

\subsection{Executables in the game tree}
\label{sec:game-tree}

The recursive glob search over the game tree returned 7 executables. They are
listed in Table~\ref{tab:game-tree}, grouped by their role.

\begin{table}[htbp]
\centering
\caption{Executables found in the game tree (7 files).}
\label{tab:game-tree}
\begin{tabular}{@{}lll@{}}
\toprule
\textbf{Executable} & \textbf{Location} & \textbf{Role} \\
\midrule
\texttt{Cyberpunk2077.exe} & \texttt{bin/x64/} & Main game executable (launcher of interest) \\
\texttt{REDprelauncher.exe} & game root & Launcher \\
\texttt{REDEngineErrorReporter.exe} & \texttt{bin/x64/} & Error reporting \\
\texttt{CrashReporter.exe} & \texttt{bin/x64/CrashReporter/} & Crash reporting \\
\texttt{7za.exe} & \texttt{bin/x64/CrashReporter/} & Archive utility \\
\texttt{redMod.exe} & \texttt{tools/redmod/bin/} & Modding tool \\
\texttt{scc.exe} & \texttt{tools/redmod/bin/} & Modding tool \\
\bottomrule
\end{tabular}
\end{table}

\subsection{Executables in the Redist folder}
\label{sec:redist}

The separate search over the \texttt{Redist} folder returned 5 executables.
These are runtime dependencies rather than game binaries. They are listed in
Table~\ref{tab:redist}.

\begin{table}[htbp]
\centering
\caption{Executables found in the \texttt{Redist} folder (5 files).}
\label{tab:redist}
\begin{tabular}{@{}lll@{}}
\toprule
\textbf{Executable} & \textbf{Location} & \textbf{Role} \\
\midrule
\texttt{vc\_redist.x64.exe} & \texttt{Redist/} & Visual C++ runtime (x64) \\
\texttt{vcredist\_x86.exe} & \texttt{Redist/} & Visual C++ runtime (x86) \\
\texttt{PhysX\_9.08.14\_9.09.0814\_SystemSoftware.exe} & \texttt{Redist/} & PhysX system software \\
\texttt{DXWebSetup.exe} & \texttt{Redist/} & DirectX web installer \\
\texttt{DXSETUP.exe} & \texttt{Redist/DirectX/} & DirectX setup \\
\bottomrule
\end{tabular}
\end{table}

\subsection{Summary}
\label{sec:summary}

In total, the audit located 12 executable files: 7 within the game tree and 5
within the \texttt{Redist} folder. The launcher of interest is
\texttt{Cyberpunk2077.exe} (located at \texttt{bin/x64/Cyberpunk2077.exe});
for the setup steps, \texttt{Run Me!.bat} is available.

\section{Discussion}
\label{sec:discussion}

The observed executable layout matches the convention of a standard
pre-installed Cyberpunk~2077 repack. Three structural observations follow.

First, the game binaries are concentrated under \texttt{bin/x64}: the main
executable \texttt{Cyberpunk2077.exe}, the error reporter, and the crash
reporter all reside there, with the launcher \texttt{REDprelauncher.exe} at the
game root. This separation of launcher and game binary is typical of
CD~Projekt~RED's distribution layout.

Second, the modding toolchain is isolated under
\texttt{tools/redmod/bin}, containing \texttt{redMod.exe} and
\texttt{scc.exe}. Their presence confirms that the repack ships the REDmod
tooling alongside the game rather than requiring a separate download.

Third, the \texttt{Redist} folder contains only runtime dependencies---Visual
C++ redistributables, the PhysX system software, and DirectX installers---and
no game code. These files support the game's execution but are not part of the
game itself, which is consistent with the role of a redistributables folder in
pre-installed distributions.

The identification of \texttt{Cyberpunk2077.exe} as the launcher of interest
answers RQ3 directly; \texttt{Run Me!.bat} serves the setup steps that precede
first launch. The primary limitation of this study is that it is a single
instance: the findings describe one installation on one volume and were
obtained by filename enumeration only, without integrity or execution checks.

\section{Conclusion}
\label{sec:conclusion}

This paper documented a recursive glob audit of executable files in an
extracted Cyberpunk~2077 installation. The audit located 12 executables in
total: 7 game-related binaries in the game tree and 5 redistributable runtime
installers in the \texttt{Redist} folder. The layout is consistent with a
standard pre-installed repack, with game binaries under \texttt{bin/x64},
modding tools under \texttt{tools/redmod/bin}, and runtime dependencies under
\texttt{Redist/}. The primary launcher of interest is
\texttt{Cyberpunk2077.exe}, with \texttt{Run Me!.bat} available for the setup
steps.

\begin{thebibliography}{9}

\bibitem{priorsession}
Prior session work (unpublished session notes): extraction of the
Cyberpunk~2077 game archive from a zstd-compressed ZIP archive to
\texttt{/Volumes/One Touch/mdata/games/cbp/Cyberpunk 2077/}.
% TODO: The source session cited no external literature. This bibliography
% contains only an internal pointer to the prior session work; any additional
% entries must be grounded in the session summary before being added.

\end{thebibliography}

\end{document}